\section{引言}
\label{sec:intro}

血吸虫病仍是流行区的重要寄生虫负担,而\emph{日本血吸虫}(\emph{Schistosoma
japonicum})感染会引起一种独特的肝损伤形式。滞留于汇管区的虫卵引发肉芽肿性
炎症与进行性间隔纤维化,形成本病在 B 型超声上以网格状、网络状及龟甲样回声
图案为特征的典型表现。由于超声检查廉价、便携且无电离辐射,它是流行区大规模
筛查与纵向随访的实用模态;在流行区内,瞬态弹性成像与磁共振成像则很少
可用~\cite{schistoelasto2018,schistopocus2024}。

然而,基于超声的分期高度依赖操作者经验。观察者间可重复性也曾在丙型肝炎
相关肝纤维化的剪切波弹性成像中得到评估:Kaposi 等研究的点剪切波方法报告了
极佳的一致性~\cite{interobserver2019};而一项针对二维剪切波弹性成像的
meta 分析报告了各研究间诊断性能存在异质性,采样协议与设备差异被讨论为可能
的影响因素~\cite{swe2018staging}。这些研究关注的是弹性成像,而非 B 型图像
的自动分级。另有一类工作将深度学习用于基于超声的纤维化评
估~\cite{usfibrosisgan2022,microflow2023,hfus2025,liverspleen2025}。

有两个特性将血吸虫分级任务与一般的纤维化分类区分开来,二者共同塑造了本
工作。

第一,临床有意义的目标是\emph{细粒度的}。常规四级量表虽可解释,但相对于
间隔纤维化的连续进展而言偏粗;以 $0.1$ 为步长量化得到的 36 级量表能更好反映
渐进变化,同时仍可映射回熟悉的临床分级。我们先前的工作 SFibAI 已证明,此类
细粒度连续评分在图像层面可行:通过预测 36 个级别上的后验分布并报告其期
望~\cite{xu2026sfibai}。恰恰因为目标是细粒度的,判别性证据也更细微:模型
必须分辨间隔纹理与实质结构的差异,而这些差异是四分类决策边界完全无需区分的。

第二,采集具有\emph{解剖结构性}。图像按标准化六部位协议采集,因此每张图像
都带有已知的解剖部位;纤维化改变以粗边界框标注。这两个信号在训练时均可得,
且都可能有信息量:正常实质的外观在不同扫查部位间存在系统性差异,因此知道
图像取自何处,会改变"什么应被视为异常"的判断。

由此引出本文研究的问题。解剖部位与病灶范围可以方便地作为\emph{训练}监督,
但不能假定它们在部署时仍然存在:常规实践中部位元数据常缺失或不可靠,框标注
则是科研阶段产物。因此,一个有用的设计必须能从这些信号中学习,同时又不依赖
它们进行推理。此类解剖监督是否真能帮助细粒度分级,以及两种信号是否以同样
方式起作用,并非先验可知。

我们提出 \textbf{SynAP-Fib}:在 SFibAI 分级骨干网络上增设两个解剖先验
分支——六类部位头与弱框病灶注意头。每个分支通过带双向梯度隔离边界的门控
残差注入向分级通路供给信息,使各分支提供解剖上下文,而不成为进入共享表征的
额外梯度路径。推理时,分级通路仅消费模型自身\emph{预测的}部位后验与病灶
注意力;不需要任何部位元数据或框标注。

将两个先验分别与联合进行评估,得到了本文的核心实证观察。相对于基线,单独
加入任一先验均使复合分级风险恶化;而同时提供二者并未继承这种恶化,并在四种
配置中取得最佳观察复合风险。因此,联合配置的行为不是两个单先验配置在孤立
情形下所能预测的。

我们的贡献如下。

\begin{itemize}
\item 我们提出 SynAP-Fib,一种细粒度超声肝纤维化分级模型,其中解剖部位与
      弱病灶证据既作为训练时监督,又经由门控、梯度隔离的残差注入作为推理时
      上下文,部署时无需任何解剖元数据。
\item 我们报告了两个解剖先验的 $2\times2$ 比较,呈现一种非可加的实证模式:
      任一先验单独使用均使复合分级风险恶化,而二者组合则不然,并取得四种
      配置中的最佳观察复合风险。
\item 我们证明联合配置以 $5.9\%$ 的参数量与 $2.8\%$ 的算术开销,在保持基线
      水平分级精度的同时,额外提供两个临床上可解释的输出——六类解剖部位与
      弱病灶定位。
\item 我们在预设的、临床加权的复合风险下进行评估,该风险综合了图像级、
      患者级与中心平衡三个视角;检查点选择仅限于验证集,测试集仅用于报告。
\end{itemize}

本文其余部分安排如下。第~\ref{sec:related} 节将 SynAP-Fib 置于已有工作
背景中。第~\ref{sec:methods} 节描述队列、架构与评估协议。第~\ref{sec:results}
节报告结果,第~\ref{sec:discussion} 节讨论其解释与局限。
